\unit{Conditional optics: how measurement becomes a gate}
\unitref{W5L3; Week 6 tutorial; Assignment 2}
\subsection{A surprising capability of linear optics}
Beam splitters and phase shifters implement
linear unitary transformations of optical
modes. They cannot directly create a
deterministic interaction between two
photons. But combining these optical
elements with ancilla photons and
\emph{postselected} photodetection can
implement an effective nonlinear operation
on a conditional subensemble.
The cost is that some runs must be rejected.

\subsection{One outcome, one Kraus operator}
If the joint state of signal and ancilla
is \(U(\ket\psi\otimes\ket{\rm anc})\)
and the ancilla detector yields \(k\),
the signal's unnormalised output is
\[
\ket{\tilde\psi_k}
={}_{\rm anc}\bra k U
(\ket\psi\otimes\ket{\rm anc}),
\]
or \(\tilde\rho_k=K_k\rho K_k^\dagger\).
The probability of this outcome is
\(p_k=\operatorname{Tr}\tilde\rho_k\);
\(\rho_k=\tilde\rho_k/p_k\) is the
normalised state for successful trials.
Do \emph{not} normalise before calculating
the success probability.

\subsection{Photon subtraction}
Consider a signal
\(\ket\psi=\sum_n c_n\ket n\)
mixed with vacuum at a weakly reflecting
beam splitter of transmissivity \(\eta\).
Conditioning on one photon in the
reflected output produces
\[
\boxed{K_{\rm click}\ket n\propto
\sqrt n\,\sqrt{1-\eta}\,
\eta^{(n-1)/2}\ket{n-1}.}
\]
The factor \(\sqrt n\) is the
annihilation-operator matrix element.
The factor \(\eta^{(n-1)/2}\) accounts for
the photons retained in the signal.
For \(\eta\) near 1, the operation
approximates \(a\ket\psi\), up to
an overall success amplitude.
It \emph{does not} remove a photon
deterministically from every trial.

\begin{worked}{Distinguishing nonorthogonal states probabilistically}
Take the two input states
\[
\ket{\psi_\pm}=\frac{2\ket0\pm\ket1}{\sqrt5},
\qquad \langle\psi_+|\psi_-\rangle=3/5.
\]
They cannot be perfectly distinguished by
a deterministic measurement. Use a
successful filter
\[
K_s=\tfrac12\ket0\bra0+\ket1\bra1.
\]
The unnormalised output is
\((\ket0\pm\ket1)/\sqrt5\)
with probability \(p_s=2/5\).
After normalising,
\[
\boxed{\ket{\phi_\pm}=
\frac{\ket0\pm\ket1}{\sqrt2}}.
\]
These are \emph{orthogonal}, so a
Hadamard-basis measurement now identifies
the sign, conditional on success.
The failure probability is \(3/5\);
the identical failure output contains
no information about the sign. Probabilistic
filtering does not violate no-cloning
or deterministic distinguishability.
\end{worked}

\subsection{Heralding and gate resources}
Conditional sign gates can transform selected
Fock coefficients with relative phases,
for example \(\alpha\ket0+\beta\ket1+
\gamma\ket2\) into
\(\alpha\ket0+\beta\ket1-\gamma\ket2\)
in a successful ancilla detection branch.
Such number-dependent relative phases
are effectively nonlinear operations,
even though the optical elements
before measurement were linear.
This idea is used in measurement-induced
photonic computation, with ancilla states,
feed-forward and suitable encoding.

\subsection{Entanglement distillation}
Two noisy Bell pairs can be processed
with bilateral gates and local measurements
to keep one pair only when parity checks
succeed. For the simple ensemble
\[
\rho=p\ket{\Phi^+}\bra{\Phi^+}
+(1-p)\ket{\Psi^+}\bra{\Psi^+},
\]
a recurrence protocol keeps matching
error labels with
\[
p_{\rm success}=p^2+(1-p)^2,
\quad
\boxed{p'=\frac{p^2}{p^2+(1-p)^2}}.
\]
For \(p=0.8\), the retained fidelity
rises to \(p'\approx0.9412\),
but the pair is obtained with probability
\(0.68\) from two initial pairs.
Entanglement quality improves at the
expense of throughput.

\begin{checkpoint}
Always report \emph{both} the conditional
output state and its success probability.
Distinguish a trace-preserving quantum
channel from a selected, trace-decreasing
instrument outcome.
\end{checkpoint}
